% !Mode:: "TeX:UTF-8" 

\BiChapter{图表公式}{Figures, Tables and Equations}

论文写作过程中最重要的是{\hei 图、表、公式}等内容的编排。

%=========================================================================================
\BiSection{图}{Figures}
\BiSubsection{单幅图}{Single Figure}

图 \ref{fig_ch1_echoes} 是用 Tikz 画的，Visio 画不出这么好看的图。
%
\begin{figure}[!ht]
	\centering
	\includegraphics[width=\textwidth]{echoes}
	\caption{图注是五号字。} \label{fig_ch1_echoes}
\end{figure}

%-----------------------------------------------------------------------------------------
\BiSubsection{多幅图}{Multiple Figures}

如果一幅图中包含多幅子图，每一幅子图都要有图注，并且子图用 (a)、(b)、(c) 等方式编号，如图 \ref{fig_ch1_badge} 所示。

\begin{figure}[!ht]
	\centering
	\subfigure[灰色的交大校徽]{\includegraphics[width=0.45\textwidth]{xjtu_gray}} \hfill
	\subfigure[蓝色的交大校徽]{\includegraphics[width=0.45\textwidth]{xjtu_blue}}
	\caption{交大校徽 \label{fig_ch1_badge}}
\end{figure}

%=========================================================================================
\BiSection{表}{Tables} 

表格要求采用三线表，与文字齐宽，顶线与底线线粗是 $1\frac12$ 磅，中线线粗是 1 磅，如表 \ref{tab_ch1} 所示\footnote{{\color{red}注意}：图表中的变量与单位通过斜线 $/$ 隔开。}。
%
\begin{table}[!ht]
	\renewcommand{\arraystretch}{1.2}
	\centering\wuhao
	\caption{表题也是五号字} \label{tab_ch1} \vspace{2mm}
	\begin{tabularx}{\textwidth}{*{4}Y}
	\toprule[1.5pt]
		Interference & DOA / degree & Bandwidth / MHz & INR / dB \\
	\midrule[1pt]
		1 & $-30$ & 20 & 60 \\
		2 & 20 & 10 & 50 \\
		3 & 40 & 5 & 40 \\
	\bottomrule[1.5pt]
	\end{tabularx}
\end{table}

%{\color{red} 通过和研究生院学位办负责论文模板的老师讨论确定：英文论文的格式和中文论文的格式完全相同，写英文论文时把 format.tex 中第 144 和 145 行的注释取消即可。}

%=========================================================================================
\BiSection{公式}{Equations}
\BiSubsection{单个公式}{Equations}

单个公式的编号如式 (\ref{equ_ch1_pdf}) 所示，该式是标准正态分布的概率密度函数\citeup{Manolakis2005}，从公式可知 \LaTeX 排版的公式比 MathType 美观，而且公式编写效率更高。
%
\begin{equation} \label{equ_ch1_pdf}
	f_Z(z) = \frac{1}{\pi\sigma^2} \exp\left(-\frac{|z-\mu|^2}{\sigma^2}\right)
\end{equation}
%
式中：$\mu$ 是 Gauss 随机变量 $Z$ 的均值；$\sigma^2$ 是 $Z$ 的方差。

%-----------------------------------------------------------------------------------------
\BiSubsection{多个公式}{Subequations}

多个公式作为一个整体可以进行二级编号，如 (\ref{equ_ch1_fourier}) 所示，该式是连续时间 Fourier 变换的正反变换公式\citeup{Vetterli2014}。
%
\begin{subequations} \label{equ_ch1_fourier}
	\begin{align}
		X(f) &= \int_{-\infty}^{\infty}x(t)e^{-j2\pi f t}\dif t \\
		x(t) &= \int_{-\infty}^{\infty}X(f)e^{j2\pi f t}\dif f
	\end{align}
\end{subequations}
%
式中：$x(t)$ 是信号的时域波形；$X(f)$ 是 $x(t)$ 的 Fourier 变换。
